\documentclass{article}

\usepackage[utf8]{inputenc}
\usepackage[german]{babel}
\usepackage{listings}
\usepackage{xcolor}

\lstset{
    basicstyle=\ttfamily\small,
    breaklines=true,
    backgroundcolor=\color{gray!10}
}

\begin{document}
    \begin{center}
        \Large \textbf{Huntsman Challenge 5: Infrastructure Lifecycle Management} \\ [6pt]
        \normalsize
        Autor: Maximilian Jakob Maag B.Sc. \\
        Version: 1.0
    \end{center}

    \section{Intro}
    Diese Challenge behandelt den gesamten Lebenszyklus von Infrastruktur - von der Bereitstellung bis zur Außerbetriebnahme.
    Sie werden lernen, wie man Infrastruktur-Änderungen systematisch verwaltet, dokumentiert und automatisiert.

    \section{Aufgaben}

    \subsection{Aufgabe 1}
    Erstellen Sie ein Server Onboarding Playbook für neue vom Rabbit deployten Server.
    \\ \\
    Das Playbook sollte automatisch:
    \begin{itemize}
        \item Den Server zum Inventory hinzufügen
        \item Basis-Konfiguration durchführen
        \item Monitoring Agent installieren
        \item Zugriff für Huntsman Team konfigurieren
        \item Initiale Compliance Checks durchführen
        \item Dokumentation aktualisieren
    \end{itemize}

    \subsection{Aufgabe 2}
    Implementieren Sie ein Change Management Playbook.
    \\ \\
    Das Playbook sollte:
    \begin{itemize}
        \item Alle Konfigurationsänderungen in einer Change Management DB registrieren
        \item Vor kritischen Änderungen ein Backup durchführen
        \item Abhängigkeiten zwischen Services prüfen
        \item Automatische Rollback Punkte erstellen
        \item Change Approval Workflow implementieren (z.B. über externe APIs)
        \item Alle Änderungen auditieren und loggen
    \end{itemize}

    \subsection{Aufgabe 3}
    Erstellen Sie ein Configuration Management Database (CMDB) Playbook.
    \\ \\
    Das System sollte folgende Daten verwalten:
    \begin{itemize}
        \item Alle verwalteten Hosts und ihre Eigenschaften
        \item Installierte Software und Versionen
        \item Konfigurationsabhängigkeiten
        \item Service Dependencies (z.B. Web Server braucht Datenbank)
        \item Historical Data über Konfigurationsänderungen
    \end{itemize}

    \subsection{Aufgabe 4}
    Implementieren Sie ein Server Decommissioning Playbook.
    \\ \\
    Das Playbook sollte:
    \begin{itemize}
        \item Abhängigkeiten zu diesem Server identifizieren
        \item Workloads auf andere Server migrieren
        \item Finale Backups erstellen
        \item Server aus dem Inventory entfernen
        \item Daten sicher löschen (wegen GDPR/Compliance)
        \item Dokumentation aktualisieren
        \item Infrastructure as Code (Terraform) aktualisieren
    \end{itemize}

    \subsection{Aufgabe 5}
    Erstellen Sie ein Infrastructure as Code (IaC) Synchronisations Playbook.
    \\ \\
    Das Playbook sollte:
    \begin{itemize}
        \item Ansible Inventory mit Terraform State synchronisieren
        \item Drift Detection durchführen (Ist der aktuelle Zustand wie dokumentiert?)
        \item Automatisch erkannte neue Server zum Inventory hinzufügen
        \item Gelöschte Server aus dem Inventory entfernen
        \item Regelmäßig Reports über Drift generieren
    \end{itemize}

    \subsection{Aufgabe 6}
    Implementieren Sie ein Service Migration Playbook.
    \\ \\
    Das Playbook sollte Szenarien unterstützen wie:
    \begin{itemize}
        \item Server upgrades ohne Ausfallzeit (Blue-Green Deployment)
        \item Regionale Migrationen (z.B. US zu EU)
        \item Applikations-Updates mit Zero Downtime
        \item Rollback Prozeduren
    \end{itemize}

    \subsection{Aufgabe 7}
    Erstellen Sie ein Capacity Planning Playbook.
    \\ \\
    Das Playbook sollte:
    \begin{itemize}
        \item Historische Ressourcennutzung analysieren
        \item Trends vorhersagen
        \item Warnungen bei Kapazitätsproblemen generieren
        \item Empfehlungen für Server-Upgrades oder Neubeschaffung geben
        \item Costs basierend auf Ressourcennutzung tracken
    \end{itemize}

    \subsection{Aufgabe 8}
    Dokumentieren Sie Ihren gesamten Infrastructure Lifecycle und erstellen Sie Standard Operating Procedures (SOPs).
    \\ \\
    Die Dokumentation sollte enthalten:
    \begin{itemize}
        \item Server Onboarding Prozess (Step-by-Step)
        \item Change Management Procedure
        \item Decommissioning Prozess
        \item Service Migration Procedure
        \item Escalation Contacts und Zeiten
        \item Häufige Probleme und deren Lösungen
    \end{itemize}

    \subsection{Aufgabe 9}
    Erstellen Sie automatisierte Tests und Validierungen für Ihre Playbooks.
    \\ \\
    Verwenden Sie Test-Tools wie:
    \begin{itemize}
        \item Molecule für Ansible Role Testing
        \item Terratest für Terraform Testing
        \item Compliance Testing Frameworks
        \item Integration Tests für komplexe Workflows
    \end{itemize}

\end{document}
